\subsection{2.4. Vulnerabilidades de seguridad en contratos inteligentes
de
Ethereum}\label{vulnerabilidades-de-seguridad-en-contratos-inteligentes-de-ethereum}

Los contratos inteligentes en Ethereum presentan un modelo de seguridad
distinto al del software tradicional. Su ejecución es pública, el código
es accesible y cualquier error puede ser analizado y explotado por
actores con incentivos económicos directos. Además, la inmutabilidad
dificulta la corrección de fallos una vez desplegado el contrato, y la
interacción con otros contratos y servicios externos incrementa la
superficie de ataque. {[}10{]}, {[}30{]}

Por este motivo, una vulnerabilidad en Solidity no debe entenderse solo
como un error de programación, sino como una debilidad explotable en un
entorno adversarial. En la práctica, los problemas de seguridad no se
limitan a fallos técnicos, sino que también incluyen errores de diseño y
de lógica de negocio. {[}31{]}

Para su análisis, resulta útil agrupar las vulnerabilidades en cuatro
categorías principales:

\begin{itemize}
\tightlist
\item
  Vulnerabilidades técnicas de ejecución
\item
  Vulnerabilidades de control y privilegios
\item
  Vulnerabilidades económicas y dependencia del entorno
\item
  Errores lógicos de negocio
\end{itemize}

\subsubsection{2.4.1. Vulnerabilidades técnicas de
ejecución}\label{vulnerabilidades-tuxe9cnicas-de-ejecuciuxf3n}

Este grupo incluye errores directamente relacionados con la ejecución en
la EVM y la implementación en Solidity.

Una de las más conocidas es la~\textbf{reentrancy}, que ocurre cuando un
contrato realiza una llamada externa antes de actualizar su estado
interno. Esto permite que el contrato receptor vuelva a ejecutar la
función vulnerable con un estado inconsistente. A pesar de que existen
patrones de mitigación bien conocidos, sigue siendo una de las
vulnerabilidades más relevantes en Ethereum. {[}32{]}

Otra categoría importante son los~\textbf{problemas aritméticos}, como
overflow y underflow. Aunque Solidity 0.8 introduce comprobaciones
automáticas, estos errores siguen siendo relevantes en contratos
antiguos, en bloques~unchecked~o en cálculos financieros incorrectos.

También destacan las~\textbf{llamadas externas inseguras}, ya que
cualquier~call~transfiere el control de ejecución a otro contrato. Esto
puede provocar comportamientos inesperados o ejecución de lógica
maliciosa. En esta línea, el uso de~delegatecall~es especialmente
crítico, ya que ejecuta código externo sobre el almacenamiento del
contrato llamador, pudiendo comprometer completamente su estado.~
{[}30{]}

Por último, los~\textbf{ataques de denegación de servicio (DoS)}~son
frecuentes en este entorno. No buscan tumbar el sistema, sino bloquear
funciones críticas, por ejemplo, mediante bucles no acotados o
reversiones provocadas por contratos externos. {[}30{]}

\subsubsection{2.4.2. Vulnerabilidades de control y
privilegios}\label{vulnerabilidades-de-control-y-privilegios}

Muchas vulnerabilidades reales no se deben a errores técnicos complejos,
sino a problemas en el control de acceso.

Esto incluye funciones sensibles sin restricciones, roles mal definidos
o privilegios excesivos. Un error típico es permitir que cualquier
usuario ejecute funciones críticas como transferencias de fondos o
cambios de configuración. {[}33{]}

Un caso especialmente problemático es el uso de~tx.origin~para
autenticación. Este valor representa el origen externo de la
transacción, no el llamador inmediato, lo que permite ataques mediante
contratos intermedios.

También son relevantes los problemas en contratos upgradeables,
especialmente los relacionados con inicializadores. Si una
función~initialize~no está correctamente protegida, un atacante puede
ejecutarla y asumir el control del contrato. {[}34{]}

\subsubsection{2.4.3. Vulnerabilidades económicas y dependencia del
entorno}\label{vulnerabilidades-econuxf3micas-y-dependencia-del-entorno}

En muchos casos, el problema no está en el código en sí, sino en cómo
interactúa el contrato con su entorno.

El ejemplo más claro es el~\textbf{front-running}, donde un atacante
observa transacciones en la mempool y envía otras con mayor prioridad
para ejecutarse antes. Este fenómeno forma parte del problema más amplio
del~\textbf{MEV (Maximal Extractable Value)}~y afecta especialmente a
aplicaciones DeFi. {[}35{]}

Otro riesgo importante es la~\textbf{manipulación de oráculos}. Si un
contrato depende de precios externos que pueden alterarse temporalmente,
un atacante puede aprovechar esa situación para obtener beneficios
indebidos, por ejemplo, en protocolos de préstamo. {[}33{]}

En esta categoría también se incluyen los problemas
de~\textbf{aleatoriedad insegura}~y el uso de~block.timestamp~como
fuente de decisiones críticas. Estas variables no son completamente
fiables ni impredecibles, por lo que pueden ser manipuladas o
anticipadas en determinados escenarios.

\subsubsection{2.4.4. Errores lógicos de
negocio}\label{errores-luxf3gicos-de-negocio}

Los errores más difíciles de detectar son los relacionados con la lógica
del contrato.

En estos casos, el código puede ser correcto desde el punto de vista
técnico, pero implementar una lógica incorrecta. Esto incluye errores en
cálculos de balances, distribución de recompensas, gestión de estados o
validación de condiciones. {[}31{]}

Este tipo de vulnerabilidades es especialmente relevante porque suele
escapar a las herramientas automáticas. Además, muchos ataques reales
combinan errores lógicos con otros factores, como manipulación de
precios o condiciones de ejecución. {[}35{]}

\subsubsection{2.4.5. Conclusión}\label{conclusiuxf3n}

En conjunto, la seguridad en contratos inteligentes no puede abordarse
únicamente mediante la detección de patrones conocidos. Aunque
vulnerabilidades como reentrancy o overflow siguen siendo relevantes,
los problemas más complejos suelen estar relacionados con el diseño del
sistema, la interacción con otros componentes y la lógica de negocio.

Esto implica que una auditoría efectiva debe analizar no solo el código,
sino también su comportamiento, sus dependencias y los incentivos que lo
rodean.
